\chapter{Metodologia}
\label{cap:metodologia}

A pesquisa adota procedimentos de análise documental e comparação técnica dos requisitos de conexão coletiva, complementados por aplicação a um projeto elétrico aprovado pela distribuidora. A unidade de análise normativa corresponde às exigências relativas a empreendimentos com múltiplas unidades consumidoras na área de concessão da Enel Ceará. O estudo de caso permite examinar as implicações das diferenças identificadas sobre soluções de engenharia documentadas, sem constituir avaliação de conformidade integral do empreendimento.

\section{Natureza e abordagem da pesquisa}
\label{sec:tipo-pesquisa}

Trata-se de pesquisa aplicada, descritiva e comparativa, fundamentada em documentos técnicos. A abordagem é predominantemente qualitativa, uma vez que o objetivo central consiste em identificar, classificar e interpretar mudanças dos requisitos normativos. Dados quantitativos são empregados nas situações em que a avaliação de repercussões sobre o projeto exige a confrontação de limites de potência, demanda, corrente, dimensões ou características nominais dos equipamentos.

A comparação das especificações e o estudo de caso constituem etapas relacionadas, porém distintas. Na primeira, identificam-se as diferenças entre os documentos de referência; na segunda, examina-se a aplicabilidade dessas diferenças às características de uma instalação projetada. A ausência de determinado requisito em uma seção não será interpretada, isoladamente, como sua eliminação, sendo necessária a verificação de remissões a outras seções, tabelas, anexos e documentos complementares.

\section{Conjunto documental}
\label{sec:documentos}

As fontes normativas principais são a especificação 0126, versão 03, de 31 de março de 2022, e a especificação 1580, versão 2, de 8 de setembro de 2025. Os documentos complementares serão consultados apenas quando necessários à interpretação dos requisitos selecionados.

A Resolução Normativa ANEEL nº 1.000/2021 e suas alterações pertinentes compõem o referencial regulatório para interpretação das condições de atendimento, das normas técnicas da distribuidora e dos períodos de transição. Também integram o conjunto documental os desenhos e anexos das especificações quando contiverem exigências que complementem o texto principal.

A aplicação prática utiliza a documentação do projeto elétrico do Edifício Palatium, empreendimento residencial localizado em Sobral, Ceará. O conjunto analisado abrange o memorial de cálculo elétrico, revisão 04, as pranchas do sistema de entrada e alimentação, o registro de responsabilidade técnica e a carta de aprovação expedida pela Enel Distribuição Ceará em 25 de outubro de 2025. Os dados técnicos do empreendimento serão apresentados na análise do estudo de caso, e não nesta caracterização metodológica. O histórico das revisões anteriores do projeto não integra o objeto da pesquisa.

\section{Procedimentos de comparação normativa}
\label{sec:procedimento}

O levantamento será realizado em etapas sucessivas. Inicialmente, serão identificados objeto, área de aplicação e elementos técnicos equivalentes nas duas especificações. Em seguida, as prescrições serão agrupadas conforme os temas: enquadramento e limites de atendimento, cálculo de demanda, formas de conexão, ramais, padrão de entrada, medição, proteção, aterramento e exigências de apresentação de projeto.

Para cada tema, serão registrados os dispositivos correspondentes em cada documento, com indicação da seção, item, tabela ou desenho, quando aplicável. A comparação será efetuada entre prescrições funcionalmente equivalentes, inclusive quando inseridas em partes distintas das especificações. Quando o documento de 2025 remeter a outra norma, serão examinadas a referência e suas condições de aplicação antes de caracterizar eventual supressão ou modificação de requisito.

As diferenças serão classificadas em seis categorias: requisito mantido, alterado, acrescentado, retirado ou substituído, reorganização textual sem efeito técnico demonstrado e situação não conclusiva. A classificação estará vinculada às evidências documentais. A existência de nota regional, condição excepcional ou exigência de estudo de viabilidade será registrada expressamente, de modo a evitar generalizações que não se apliquem ao Ceará.

\section{Critérios de avaliação dos impactos técnicos}
\label{sec:impactos}

Os impactos serão avaliados segundo a possibilidade de interferência em decisões de engenharia documentadas. Serão considerados a modalidade de alimentação, a localização e a configuração do padrão de entrada, o dimensionamento de condutores e dispositivos de proteção, a disposição dos centros de medição e os requisitos de documentação e aprovação.

A classificação de determinada exigência como alterada não implicará, por si só, a necessidade de reformular o projeto utilizado como referência. Para estabelecer uma consequência técnica, será necessário identificar a solução registrada na documentação, o requisito correspondente na especificação 1580 e as condições de aplicação. Os resultados serão classificados como solução compatível com o requisito examinado, adequação potencialmente necessária, situação condicionada à avaliação da distribuidora ou informação documental insuficiente.

Quando disponíveis os parâmetros necessários, serão realizados cálculos demonstrativos com identificação das grandezas, hipóteses e fontes empregadas. Esses cálculos terão finalidade analítica e não substituirão os memoriais nem a responsabilidade técnica dos profissionais responsáveis pelo projeto original.

\section{Aplicação ao estudo de caso}
\label{sec:estudocaso}

A aplicação prática será conduzida sobre as soluções de entrada coletiva representadas no projeto aprovado do Edifício Palatium. A documentação existente será adotada como configuração de referência para verificar os possíveis efeitos das exigências selecionadas da especificação 1580 sobre um mesmo empreendimento. A especificação 0126 será considerada como referência normativa anterior, em conformidade com a delimitação da pesquisa.

A avaliação ficará restrita aos elementos presentes no memorial e nas pranchas do projeto. Para cada tema identificado como relevante na comparação normativa, será estabelecida uma relação entre a solução projetada, o requisito correspondente e a consequência técnica potencial. A pesquisa não promoverá alterações nos documentos originais nem atribuirá desconformidade ao projeto sem a confirmação das condições normativas pertinentes.

A carta de aprovação da concessionária comprova o deferimento do projeto apresentado, mas não discrimina os dispositivos normativos empregados na análise. O contexto de transição será examinado à luz das datas documentadas e do art. 20 da Resolução Normativa ANEEL nº 1.000/2021. Não serão presumidas a data efetiva de comunicação da especificação 1580 ou a incidência de hipóteses excepcionais sem suporte documental.

\section{Tratamento dos resultados e limitações}
\label{sec:limitacoes}

Os resultados da comparação serão sistematizados em quadros contendo os temas examinados, os dispositivos correspondentes das especificações 0126 e 1580, a natureza da alteração e a repercussão técnica identificada. Os resultados do estudo de caso serão apresentados separadamente, com indicação da solução documentada e das condições necessárias para cada avaliação.

A pesquisa limita-se às versões documentais selecionadas e aos critérios aplicáveis à área de concessão do Ceará. Não serão investigados os requisitos exclusivos de outras distribuidoras, o histórico de revisões internas do projeto, as instalações efetivamente executadas, os procedimentos de vistoria ou a energização do empreendimento. A análise não constitui laudo de conformidade nem substitui a avaliação técnica da concessionária em situações sujeitas a estudos de viabilidade.

Quando não houver elementos suficientes para determinada conclusão, a limitação será registrada no resultado correspondente. Esse procedimento preserva a distinção entre as informações expressamente documentadas, a interpretação de requisitos técnicos e as hipóteses dependentes de confirmação adicional.
